\chapter{Entiers, divisibilité et fractions}
\section*{Exercice 2.1 -- Diviseurs}
Les diviseurs de $36$ sont $1,2,3,4,6,9,12,18,36$. On a $84=6\times14=7\times12$ ; $84=8\times10+4$ n'est pas divisible par $8$.

\section*{Exercice 2.2 -- Dernier multiple}
$250=13\times19+3$ ; le multiple recherché est $247$. Le suivant est $260>250$, et les multiples de $13$ sont croissants avec leur quotient entier.

\section*{Exercice 2.3 -- Fraction}
$630=2\times3^2\times5\times7$ et $462=2\times3\times7\times11$. Le PGCD vaut $42$, d'où $630/462=15/11$. Les entiers $15$ et $11$ sont premiers entre eux.

\section*{Exercice 2.4 -- Parité}
Pour $2k+1$ et $2\ell+1$, la somme vaut $2(k+\ell+1)$ ; le produit vaut $4k\ell+2k+2\ell+1=2(2k\ell+k+\ell)+1$. Les quantités entre parenthèses sont entières.

\section*{Exercice 2.5 -- Lots}
Le nombre de lots divise les deux effectifs ; son maximum est $\PGCD(96,144)=48$. Chaque lot contient $96/48=2$ rouges et $144/48=3$ bleues.

\section*{Exercice 2.6 -- Un reste invariant}
Écrire $n=2k+1$ donne $n^2-1=4k(k+1)$. Parmi $k$ et $k+1$, l'un est pair, donc $k(k+1)=2m$ pour un entier $m$. Ainsi $n^2-1=8m$.
